% Theme: Notebook  --  best for worked exercises, problem + solution
% Build:  ./build.sh example-notebook.tex   (or compile on Overleaf)
\documentclass[12pt]{beamer}
\usetheme{ignotebook}

\ighandle{@quantum_danesh}
\igseries{Ejercicios resueltos}
\igissue{04}

\begin{document}

\coverslide{Ejercicio resuelto}
  {¿Es el amortiguamiento de amplitud un canal?}
  {Completitud, acción sobre $\rho$ y punto fijo, paso a paso.}

\begin{frame}{El enunciado}{Ejercicio 2.4}
Dados los operadores de Kraus
\[
  K_0 = \begin{pmatrix} 1 & 0 \\ 0 & \sqrt{1-\gamma} \end{pmatrix},
  \qquad
  K_1 = \begin{pmatrix} 0 & \sqrt{\gamma} \\ 0 & 0 \end{pmatrix},
\]
con $0 < \gamma \le 1$:
\begin{enumerate}[(a)]
  \item Comprueba que $\sum_i K_i^\dagger K_i = \id$.
  \item Calcula $\mathcal{A}_\gamma(\rho)$ para un $\rho$ general.
  \item Encuentra su punto fijo.
\end{enumerate}
\end{frame}

\begin{frame}{Completitud}{Apartado (a)}
\[
  K_0^\dagger K_0 = \begin{pmatrix} 1 & 0 \\ 0 & 1-\gamma \end{pmatrix},
  \qquad
  K_1^\dagger K_1 = \begin{pmatrix} 0 & 0 \\ 0 & \gamma \end{pmatrix}
\]
\[
  K_0^\dagger K_0 + K_1^\dagger K_1 = \mhl{\id}
\]
Así que preserva la traza. \aside{\checkmark\ correcto}
\begin{notebox}[Ojo]
No olvides la daga:
$K_1\rho K_1^\dagger \neq K_1^\dagger\rho K_1$.
\end{notebox}
\end{frame}

\begin{frame}{Acción sobre $\rho$}{Apartado (b)}
Con $\rho = \begin{pmatrix} a & b \\ b^* & 1-a \end{pmatrix}$:
\begin{align*}
  K_0\rho K_0^\dagger &= \begin{pmatrix} a & \sqrt{1-\gamma}\,b \\ \sqrt{1-\gamma}\,b^* & (1-\gamma)(1-a) \end{pmatrix}\\
  K_1\rho K_1^\dagger &= \begin{pmatrix} \gamma(1-a) & 0 \\ 0 & 0 \end{pmatrix}
\end{align*}
\uncover<2->{Sumando:
\keyeq{\mathcal{A}_\gamma(\rho) = \begin{pmatrix} a + \mhl{\gamma(1-a)} & \mhl*{\sqrt{1-\gamma}}\,b \\
  \mhl*{\sqrt{1-\gamma}}\,b^* & (1-\gamma)(1-a) \end{pmatrix}}}
\end{frame}

\begin{frame}{Punto fijo}{Apartado (c)}
Imponemos $\mathcal{A}_\gamma(\rho^\star) = \rho^\star$:
\begin{gather*}
  \sqrt{1-\gamma}\,b = b \;\Rightarrow\; b = 0,\\
  a + \gamma(1-a) = a \;\Rightarrow\; a = 1 .
\end{gather*}
\keyeq{\rho^\star = \ketbra{0}}
Iterando, $1-a_n = (1-\gamma)^n(1-a_0)$ y $b_n = (1-\gamma)^{n/2}\, b_0$:
la población \hl{cae al estado fundamental} y las coherencias se apagan.
\end{frame}

\begin{frame}{Errores típicos}{Repaso}
\begin{itemize}
  \item Confundir $K_1\rho K_1^\dagger$ con $K_1^\dagger\rho K_1$.
  \item Olvidar que $\gamma \in [0,1]$ es una probabilidad.
  \item Creer que el punto fijo depende de $\rho_0$: aquí no.
\end{itemize}
\begin{thmbox}[Comprobación final]
\(
  \Tr\mathcal{A}_\gamma(\rho) = a + \gamma(1-a) + (1-\gamma)(1-a) = 1
\)
\end{thmbox}
\end{frame}

\closingslide{¿Y si lo mezclamos?}
  {Combinado con el canal despolarizante, el punto fijo puede quedar en
   cualquier punto entre el centro y el polo norte. Lo cuento en el blog.}

\end{document}
